\section{Operator Fundamental Logika Dasar (Negasi, Konjungsi, Disjungsi)}

Proposisi atomik dapat digabungkan membentuk \textbf{proposisi majemuk} dengan menggunakan \textbf{operator penghubung logika} (*logical connectives*) \cite{rosen2019}. Tiga operator dasar yang paling sering dipakai adalah negasi, konjungsi, dan disjungsi.

\subsection{Negasi (NOT / $\neg$ / $\sim$)}

Negasi adalah operator uner: ia hanya beroperasi pada satu proposisi dan menghasilkan nilai kebenaran yang berlawanan.

Dengan simbol $(\neg)$, negasi dari $p$ ditulis $\neg p$ (dibaca: ``bukan $p$'' atau ``tidak $p$'').

\begin{definisi}[Nilai kebenaran negasi]
Jika $p$ bernilai benar (T), maka $\neg p$ bernilai salah (F). Jika $p$ bernilai salah (F), maka $\neg p$ bernilai benar (T).
\end{definisi}

\begin{contoh}[Negasi dalam bahasa alami]
Misalkan $p$ : ``Saya suka apel hijau.''\\
Negasi yang setara secara logika bukan sekadar menempelkan kata ``tidak'' secara sembarang, melainkan menyangkal klaim secara utuh, misalnya: ``Tidak benar bahwa saya suka apel hijau'' atau ``Pernyataan bahwa saya suka apel hijau tidak benar.''
\end{contoh}

\subsection{Konjungsi (AND / $\wedge$)}

Konjungsi menggabungkan dua proposisi $p$ dan $q$ menjadi $p \wedge q$ (dibaca: ``$p$ dan $q$'').

\begin{definisi}[Nilai kebenaran konjungsi]
Proposisi $p \wedge q$ bernilai benar (T) \textbf{hanya jika} $p$ dan $q$ keduanya benar (T). Jika salah satu saja bernilai salah (F), maka $p \wedge q$ bernilai salah (F).
\end{definisi}

Dalam penerapan, konjungsi sering muncul ketika suatu sistem mensyaratkan beberapa kondisi sekaligus, misalnya filter pencarian yang mengharuskan dua kriteria terpenuhi bersamaan \cite{wiki_first_order_logic}.

\subsection{Disjungsi (OR / $\vee$)}

Disjungsi menggabungkan dua proposisi $p$ dan $q$ menjadi $p \vee q$ (dibaca: ``$p$ atau $q$''). Dalam logika proposisional standar, yang digunakan adalah \textbf{disjungsi inklusif}.

\begin{definisi}[Nilai kebenaran disjungsi inklusif]
Proposisi $p \vee q$ bernilai benar (T) jika \textbf{minimal salah satu} dari $p$ atau $q$ bernilai benar (T). Proposisi $p \vee q$ baru bernilai salah (F) jika $p$ dan $q$ keduanya salah (F).
\end{definisi}

\begin{contoh}[Disjungsi dalam syarat administrasi]
Misalkan aturan menyatakan: ``Pemohon kredit harus menyerahkan BPKB kendaraan \textbf{atau} bukti kepemilikan rumah yang memenuhi persyaratan.''\\
Jika pemohon hanya memenuhi salah satu dari dua dokumen tersebut, syarat disjungsi dapat terpenuhi; jika keduanya tidak terpenuhi, syarat tidak terpenuhi.
\end{contoh}
